%! AP Statistics — Unit 4, Lesson 4.11 — Guided Notes
\documentclass[10pt]{article}
\usepackage{apstats-article}
\usepackage{apstats-boxes}

\begin{document}

\pageheader{Unit 4, Lesson 4.11}{Guided Notes}
\namedateperiod

\vspace{0.12in}

\begin{objectivebox}
\small By the end of this lesson, I will be able to\dots\\[4pt]
\writeline\\[1pt]
\writeline\\[1pt]
\writeline
\end{objectivebox}

\vspace{0.1in}

\begin{vocabbox}
\termblanklong{Mean of a binomial distribution}
\termblanklong{Standard deviation of a binomial distribution}
\termblanklong{Cumulative binomial probability}
\termblanklong{binomcdf(n, p, k)}
\end{vocabbox}

\vspace{0.1in}

\begin{hookbox}
\small\textbf{Opening scenario:} A professional archer hits the target 90\% of the time.
She shoots 50 arrows in a competition.

\vspace{6pt}
How many hits do you expect her to average across many competitions? \blank{2cm}\\[6pt]
How did you reason that? \writeline\\[1pt]\writeline
\end{hookbox}

\vspace{0.1in}

\begin{notesbox}{Mean and Standard Deviation of a Binomial Distribution}
\small
For $X \sim B(n, p)$:

\vspace{6pt}
\renewcommand{\arraystretch}{2.0}
\begin{tabularx}{\linewidth}{@{} >{\bfseries}p{0.26\linewidth} X @{}}
Mean & $\mu_X = \blank{3cm}$ \\
Standard deviation & $\sigma_X = \blank{5cm}$ \\
\end{tabularx}

\vspace{8pt}
\textbf{Archer example:} $n = 50$, $p = 0.90$
\[
  \mu_X = \blank{3cm} = \blank{2cm}
  \qquad
  \sigma_X = \sqrt{\blank{3cm}} = \blank{2cm}
\]

\vspace{6pt}
\textbf{Interpret $\mu_X$ in context:}\\[4pt]
\writeline\\[1pt]\writeline

\vspace{6pt}
\textbf{Interpret $\sigma_X$ in context:}\\[4pt]
\writeline
\end{notesbox}

\vspace{0.1in}

\begin{notesbox}{Cumulative Binomial Probabilities}
\small

\textbf{Problem:} What if we want $P(X \ge 48)$, not just $P(X = 48)$?

\vspace{6pt}
Adding individual probabilities is impractical. Use the calculator.

\vspace{8pt}
\renewcommand{\arraystretch}{1.8}
\begin{tabularx}{\linewidth}{@{} >{\bfseries}p{0.35\linewidth} X @{}}
\texttt{binomcdf(n, p, k)} & gives $P(X \;\blank{1.5cm}\; k) = $ sum from $0$ to $k$ \\
$P(X \ge k)$ & $= 1 - P(X \le \blank{1.5cm}) = 1 - \texttt{binomcdf(n, p,} \blank{1cm}\texttt{)}$ \\
$P(X > k)$ & $= 1 - P(X \le \blank{1.5cm}) = 1 - \texttt{binomcdf(n, p,} \blank{1cm}\texttt{)}$ \\
\end{tabularx}

\vspace{8pt}
\textbf{Key boundary warning:} $P(X \ge k)$ and $P(X > k)$ use different values
inside \texttt{binomcdf}. Always rewrite the inequality first.

\vspace{8pt}
\textbf{Example:} Find $P(X \ge 48)$ for the archer ($n=50$, $p=0.90$).

\vspace{4pt}
$P(X \ge 48) = 1 - P(X \le \blank{1.5cm}) = 1 - \texttt{binomcdf(}\blank{1cm}\texttt{,}\;\blank{1cm}\texttt{,}\;\blank{1cm}\texttt{)} = \blank{3cm}$
\end{notesbox}

\vspace{0.1in}

\begin{practicebox}
\small
\begin{enumerate}[leftmargin=1.5em, itemsep=10pt]

  \item A manufacturer claims 15\% of its items are returned.
        A store buys 60 items. Let $X =$ the number returned. Assume $X \sim B(60, 0.15)$.
        \begin{enumerate}[label=(\alph*), itemsep=6pt]
          \item $\mu_X = $ \blank{3cm} \quad $\sigma_X = $ \blank{4cm}
          \item Interpret $\mu_X$: \writeline\\[1pt]\writeline
          \item Find $P(X \le 5)$: \texttt{binomcdf(}\blank{1cm}\texttt{,}\;\blank{1cm}\texttt{,}\;5\texttt{)} $\approx$ \blank{2cm}
          \item Find $P(X \ge 12)$: $1 - \texttt{binomcdf(}\blank{1cm}\texttt{,}\;\blank{1cm}\texttt{,}\;\blank{1cm}\texttt{)} \approx$ \blank{2cm}
        \end{enumerate}

\end{enumerate}
\end{practicebox}

\end{document}
